\documentclass[11pt,letterpaper]{article}

\usepackage[utf8]{inputenc}
\usepackage[T1]{fontenc}
\usepackage[spanish]{babel}
\usepackage[letterpaper,margin=1in]{geometry}
\usepackage[hidelinks]{hyperref}
\usepackage{xcolor}
\usepackage{enumitem}
\usepackage{longtable}
\usepackage{booktabs}
\usepackage{fancyhdr}
\usepackage{tcolorbox}

\usepackage{etoolbox}

\pretocmd{\section}{\clearpage}{}{}

\pagestyle{fancy}
\fancyhf{}
\fancyhead[L]{Manual -- Proyecto visitas al CMAT}
\fancyhead[R]{\thepage}
\renewcommand{\headrulewidth}{0pt}
\setlength{\headheight}{14pt}

\setlength{\parindent}{0pt}
\setlength{\parskip}{0.55em}
\setlist[itemize]{itemsep=0.25em, topsep=0.25em}
\setlist[enumerate]{itemsep=0.35em, topsep=0.25em}

\definecolor{cmatblue}{HTML}{1F4E79}
\definecolor{cmatgray}{HTML}{F4F6F8}
\definecolor{cmatred}{HTML}{8A1C1C}

% greens
\definecolor{green}{HTML}{2E7D32}
\definecolor{green_light}{HTML}{C8E6C9}

\newtcolorbox{nota}{
  colback=cmatgray,
  colframe=cmatblue,
  boxrule=0.5pt,
  arc=1mm,
  left=8pt,
  right=8pt,
  top=6pt,
  bottom=6pt
}

\newtcolorbox{advertencia}{
  colback=red!4,
  colframe=cmatred,
  boxrule=0.5pt,
  arc=1mm,
  left=8pt,
  right=8pt,
  top=6pt,
  bottom=6pt
}

\newtcolorbox{summary}{
  colback=green_light,
  colframe=green,
  boxrule=0.5pt,
  arc=1mm,
  left=8pt,
  right=8pt,
  top=6pt,
  bottom=6pt
}

\newcommand{\code}[1]{\texttt{#1}}

\begin{document}

\pagestyle{empty}
\begin{center}
  \vspace*{5cm}
  {\LARGE \textbf{Manual}}\\[0.4em]
  {\large Proyecto de actualizacion de figuras del reporte CMAT}\\[0.8em]
  Proyecto: \code{proyecto visitas}
\end{center}

\tableofcontents
\pagestyle{fancy}
\newpage

\section{Como actualizar las figuras}

El proceso esta automatizado para actualizar las  figuras y compilar el reporte con un solo comando. 

Los nuevos Excel de entrada deben estar en:

\begin{verbatim}
data/
|- Asesorias2024.xlsx
`- Materias estudiantes-profesores 2019-2025 P y O.xlsx
\end{verbatim}


\begin{summary}
Para actualizar las figuras se deben de mover los nuevos datos a la carpeta \code{data/} y luego se debe ejecutar:

\begin{verbatim}
run_analysis.bat
\end{verbatim}
\end{summary}




Ese archivo esta en la raiz del proyecto. Al ejecutarlo, el proceso hace todo lo
necesario:

\begin{enumerate}
  \item Verifica si Miniconda esta instalado.
  \item Si Miniconda no existe, ejecuta \code{scripts/install/miniconda.bat}.
  \item Verifica si existe el ambiente Conda \code{cmat314}.
  \item Si el ambiente no existe, ejecuta \code{scripts/install/entorno.bat}.
  \item Activa \code{cmat314}.
  \item Instala el proyecto en modo editable con \code{python -m pip install -e .}.
  \item Ejecuta \code{python src/run\_analysis.py}.
  \item Intenta compilar \code{reporte/reporte.tex} con \code{lualatex}.
\end{enumerate}

\begin{advertencia}
Si \code{lualatex} no esta instalado, las figuras se actualizan de todos modos,
pero el PDF del reporte no se compila. En ese caso hay que instalar MiKTeX:

\begin{verbatim}
https://miktex.org/download

El reporte esta en reporte/reporte.tex, asi que se puede abrir ese archivo con
MiKTeX y compilarlo ahi.
\end{verbatim}
\end{advertencia}



\section{Que esta pasando}

El proyecto toma dos archivos Excel, genera tablas y figuras en Python, y luego
compila un reporte LaTeX que usa esos archivos generados.

\begin{summary}
Los Excel de entrada deben estar en:

\begin{verbatim}
data/
|- Asesorias2024.xlsx
`- Materias estudiantes-profesores 2019-2025 P y O.xlsx
\end{verbatim}
\end{summary}

Las salidas se escriben en:

\begin{verbatim}
outputs/
|- figures/
|- tables/
|- tex/
|- logs/
|- README.md
`- summary.json
\end{verbatim}

\begin{summary}
El reporte esta en:

\begin{verbatim}
reporte/reporte.tex
\end{verbatim}
\end{summary}

El reporte lee las figuras desde \code{../outputs/figures/} y fragmentos de
tablas desde \code{../outputs/tex/}. Por eso, despues de correr
\code{run\_analysis.bat}, \code{reporte/reporte.tex} puede compilarse sin editar rutas.

\section{Estructura del proyecto}

La estructura actual es:

\begin{verbatim}
proyecto visitas/
|- config/
|  `- settings.py
|- data/
|- fonts/
|- manual/
|  `- manual.tex
|- outputs/
|- reporte/
|  `- main.tex
|- scripts/
|  `- install/
|     |- miniconda.bat
|     `- entorno.bat
|- src/
|  |- run_analysis.py
|  `- visitas_analysis/
|- pyproject.toml
|- requirements.txt
`- run_analysis.bat
\end{verbatim}

\subsection{Archivos de la raiz}

\begin{longtable}{p{0.28\textwidth} p{0.62\textwidth}}
\toprule
\textbf{Archivo} & \textbf{Funcion} \\
\midrule
\endhead
\code{run\_analysis.bat} & Entrada principal para actualizar figuras y compilar el reporte. \\
\code{pyproject.toml} & Define el paquete Python, la version de Python y las dependencias fijadas. \\
\code{requirements.txt} & Lista practica de dependencias fijadas. \\
\bottomrule
\end{longtable}

\subsection{Scripts de instalacion}

\begin{longtable}{p{0.34\textwidth} p{0.56\textwidth}}
\toprule
\textbf{Archivo} & \textbf{Funcion} \\
\midrule
\endhead
\code{scripts/install/miniconda.bat} & Instala Miniconda en \code{\%USERPROFILE\%\textbackslash miniconda3} si no existe. \\
\code{scripts/install/entorno.bat} & Crea o repara el ambiente \code{cmat314} con Python 3.14.4 e instala el proyecto. \\
\bottomrule
\end{longtable}

\section{Python y dependencias}

El proyecto sigue una estructura estandar de paquete Python con
\code{pyproject.toml}. Eso permite instalarlo con:

\begin{verbatim}
python -m pip install -e .
\end{verbatim}

La instalacion editable significa que Python usa directamente el codigo en la
carpeta del proyecto. Si se modifica un archivo dentro de \code{src/}, no hace
falta reconstruir un paquete para probar el cambio.

La version de Python esta fijada en \code{pyproject.toml}:

\begin{verbatim}
requires-python = ">=3.14,<3.15"
\end{verbatim}


Las dependencias principales tambien estan fijadas para evitar cambios
incompatibles en el futuro:

\begin{verbatim}
pandas==3.0.3
numpy==2.4.6
matplotlib==3.10.9
seaborn==0.13.2
scipy==1.17.1
scikit-learn==1.8.0
statsmodels==0.14.6
openpyxl==3.1.5
\end{verbatim}

\begin{summary}
  Esto ya esta automatizado en \code{run\_analysis.bat}. Para entrar a tu entorno de conda y modificar el proyecto, haz:
  \begin{verbatim}
  conda activate cmat314
  \end{verbatim}
\end{summary}

\section{Entrada Python principal}

El proyecto esta organizado para que el analisis completo se ejecute desde un solo archivo Python.

\begin{summary}
Dentro de \code{src}, el archivo mas importante para ejecutar el analisis es:

\begin{verbatim}
src/run_analysis.py
\end{verbatim}
\end{summary}

Este archivo:

\begin{enumerate}
  \item Detecta la raiz del proyecto.
  \item Carga la configuracion desde \code{config/settings.py}.
  \item Llama a \code{run\_visitas\_pipeline}.
  \item Muestra en consola donde quedaron las salidas.
\end{enumerate}

Tambien acepta opciones:

\begin{verbatim}
python src/run_analysis.py --skip-raw-figures
python src/run_analysis.py --skip-report-assets
\end{verbatim}

\code{run\_analysis.bat} llama a este archivo sin opciones para
regenerar todo.

\section{Modulos principales de \code{src}}

\begin{summary}
El paquete vive en:

\begin{verbatim}
src/visitas_analysis/
\end{verbatim}
\end{summary}

\begin{longtable}{p{0.35\textwidth} p{0.55\textwidth}}
\toprule
\textbf{Modulo} & \textbf{Responsabilidad} \\
\midrule
\endhead
\code{pipeline/main\_pipeline.py} & Orquesta la corrida completa del analisis. \\
\code{paths.py} & Define rutas internas y convenciones de carpetas. \\
\code{release\_figures.py} & Define nombres estables de figuras usados por Python y LaTeX. \\
\code{io/concentrado\_reader.py} & Lee el Excel de materias de forma robusta. \\
\code{reporting/descriptive\_pipeline.py} & Construye los datos descriptivos del reporte. \\
\code{reporting/metrics.py} & Calcula metricas, conteos y tablas. \\
\code{reporting/plots.py} & Genera figuras descriptivas en PDF. \\
\code{reporting/render.py} & Escribe CSV, JSON y fragmentos \code{.tex}. \\
\code{reporting/report\_assets.py} & Genera los activos del reporte en \code{outputs/}. \\
\code{reporting/run\_log.py} & Escribe logs con fecha, estado y archivos Excel usados. \\
\code{visualization/style.py} & Registra la fuente local y aplica estilo de Matplotlib. \\
\code{analysis/report\_compatible/} & Genera figuras analiticas compatibles con el reporte. \\
\bottomrule
\end{longtable}

\section{Figuras y nombres}

\begin{summary}
Las figuras se guardan en PDF dentro de:

\begin{verbatim}
outputs/figures/
\end{verbatim}
\end{summary}

Las figuras descriptivas iniciales usan estos nombres:

\begin{verbatim}
02_01_visits_histogram.pdf
02_02_visits_histogram_low_counts.pdf
02_03_visits_ecdf.pdf
02_04_visits_tail_curve.pdf
02_05_visits_continuation_curve.pdf
02_06_visits_lorenz_curve.pdf
02_07_visits_by_year.pdf
02_08_classroom_size_distribution.pdf
\end{verbatim}

Los nombres se mantienen en:

\begin{verbatim}
src/visitas_analysis/release_figures.py
\end{verbatim}

\section{Fuente del reporte}

Para el reporte se usa la fuente de EB Garamond, no esta instalada por default en windows, esta dentro del proyecto:

\begin{verbatim}
fonts/static/
\end{verbatim}

LaTeX y Matplotlib usan esa fuente local para no depender de que Windows la tenga
instalada.



\section{Mantenimiento}


\textbf{Para cambiar una figura:}

\begin{enumerate}
  \item Ubicar la funcion que la genera en \code{src/visitas\_analysis/}.
  \item Mantener la salida en PDF, para el reporte.
  \item Si cambia el nombre, actualizar \code{release\_figures.py}.
  \item Actualizar \code{reporte/reporte.tex} solo si el nombre usado por LaTeX cambia.
  \item Ejecutar \code{run\_analysis.bat}.
\end{enumerate}


\textbf{Para cambiar dependencias:}

\begin{enumerate}
  \item Editar \code{pyproject.toml}.
  \item Mantener \code{requirements.txt} sincronizado.
  \item Probar con \code{run\_analysis.bat}.
\end{enumerate}

\begin{summary}
Los datos entran por \code{data/}, todo lo generado sale
por \code{outputs/}, y el reporte compila desde \code{reporte/reporte.tex}.
\end{summary}

\end{document}
